\section{Preliminaries}

Throughout the section, we let $R$ be a ring, and let $f,g\in R[x]$ be two non-zero univariate polynomials of degree $p$ and $q$, respectively. Thus, we can write $f=a_p x^p + a_{p-1} x^{p-1} + \cdots + a_0$ and $g=b_q x^q + b_{q-1} x^{q-1} + \cdots + b_0$, with $a_p, b_q \neq 0$.

\subsection{Subresultant polynomials}

\begin{definition}[Sylvester matrix]
    The Sylvester matrix associated with $f$ and $g$ is the $(p+q)\times(p+q)$ matrix defined as follows:
    \begin{equation*}
        \Syl(f,g)=
        \begin{pmatrix}
            a_p     &        &        &        & b_q     &        &        &        \\
            a_{p-1} & a_p    &        &        & b_{q-1} & b_q    &        &        \\
            \vdots  & \vdots & \ddots &        & \vdots  & \vdots & \ddots &        \\
            a_0     & a_1    & \cdots & a_p    & b_0     & b_1    & \cdots & b_q    \\
                    & a_0    & \ddots & \vdots &         & b_0    & \ddots & \vdots \\
                    &        & \ddots & a_1    &         &        & \ddots & b_1    \\
                    &        &        & a_0    &         &        &        & b_0
        \end{pmatrix},
    \end{equation*}
    where all unspecified entries are zero.
\end{definition}

In Sylvester's original paper \cite{sylvester1853xviii}, no explicit row or column convention was specified for the Sylvester matrix. Consequently, some sources in the literature (e.g., \cite{basu2006algorithms}) present the Sylvester matrix as the transpose of the form given above. In this report, we adopt a column-oriented convention, under which the Sylvester matrix is the matrix of the linear map
\begin{align*}
    \phi : R[x]_{<q} \oplus R[x]_{<p} & \to R[x]_{<p+q} \\
    (u,v)                             & \mapsto uf + vg
\end{align*}
in the canonical monomial basis $(1,x,x^2,\ldots)$. This leads to the well-known definition of the resultant as the determinant of the Sylvester matrix.

\begin{definition}[Resultant]
    The resultant of $f$ and $g$, denoted by $\Res(f,g)$, is defined as the determinant of the Sylvester matrix $\Syl(f,g)$.
\end{definition}

As for the subresultant polynomials, they are defined using the determinants of certain submatrices of the Sylvester matrix, called the \emph{Sylvester-Habicht matrices}, which are defined as follows.

\begin{definition}[Sylvester-Habicht matrix]
    Given $0\leq i\leq q$, the $i$-th Sylvester-Habicht matrix $S_i(f,g)$ is the $(p+q-i)\times(p+q-2i)$ submatrix of $\Syl(f,g)$ obtained as follows: take the first $p-i$ columns of $\Syl(f,g)$ without their last $i$ zero rows, then append the next $q-i$ columns starting from the $(p+1)$-th row of $\Syl(f,g)$ without their last $i$ zero rows.
\end{definition}

For $0\leq j\leq i$, we denote by $M_{i,j}(f,g)$ the $(p+q-2i)\times(p+q-2i)$ submatrix of $S_i(f,g)$ obtained by taking the first $p+q-2i-1$ rows of $S_i(f,g)$ and the $(p+q-i-j)$-th row of $S_i(f,g)$. Then, $M_{i,i}(f,g)$ is the matrix of the linear map
\begin{align*}
    \phi_i : R[x]_{<q-i} \oplus R[x]_{<p-i} & \to R[x]_{<p+q-2i} \\
    (u,v)                                   & \mapsto uf + vg
\end{align*}
in the canonical monomial basis. In the same spirit as for resultants, the subresultants and subresultant polynomials of $f$ and $g$ are defined using the determinants of $M_{i,j}(f,g)$.

\begin{definition}[Subresultant]
    Given $0\leq i\leq q$, the $i$-th subresultant of $f$ and $g$, denoted by $\sRes_i(f,g)$, is defined as the determinant of the matrix $M_{i,i}(f,g)$.
\end{definition}

\begin{definition}[Subresultant polynomial]
    Given $0\leq i\leq q$, the $i$-th subresultant polynomial of $f$ and $g$, denoted by $\sResP_i(f,g)$, is defined as
    \begin{equation*}
        \sResP_i(f,g) = \sum_{j=0}^i \det(M_{i,j}(f,g)) x^j.
    \end{equation*}

    We say that $\sResP_i(f,g)$ is non-defective if $\deg(\sResP_i(f,g))=i$ (i.e., $\sRes_i(f,g)\neq 0$), and defective otherwise.
\end{definition}

\subsection{Structure theorem for subresultants}

While subresultants are defined as the determinants of certain submatrices of the Sylvester matrix, they also have a nice structure that relates them to the GCD and the remainders in the Euclidean algorithm, allowing us to compute them iteratively without performing any division in the fraction field of $R$. This is known as the structure theorem for subresultants, which is stated as follows.

\begin{theorem}[{\cite[Thm. 8.56]{basu2006algorithms}}]
    Let $0\leq i\leq q$. Suppose that $\sResP_{i-1}(f,g)$ is non-zero and of degree $0\leq j\leq i-1$. Then,
    \begin{itemize}
        \item If $\sResP_{j-1}(f,g)$ is zero, then $\sResP_{i-1}(f,g)\propto\gcd(f,g)$.

              Moreover, $\sResP_{k}(f,g)=0$ for all $k\leq j-1$.

        \item If $\sResP_{j-1}(f,g)$ is non-zero and of degree $0\leq k\leq j-1$, then $\sResP_{k-1}(f,g)$ is proportional to $\Rem(\sResP_{i-1}(f,g),\sResP_{j-1}(f,g))$.

              Moreover, if $\sResP_{j-1}(f,g)$ is defective, then $\sResP_{k}(f,g)\propto\sResP_{j-1}(f,g)$, and for all $k+1\leq\ell\leq j-2$, $\sResP_{\ell}(f,g)=0$.
    \end{itemize}
\end{theorem}

An immediate consequence is that non-defective subresultant polynomials $\sResP_{k-1}(f,g)$, as well as the $\sResP_{k}(f,g)$, are essentially the remainders in the Euclidean algorithm, up to a scaling factor. The latter sequence of subresultant polynomials, shown in Example \ref{ex:subres}, is known as the \emph{subresultant polynomial remainder sequence} of $f$ and $g$.

We refer the reader to \cite[Chapter 8]{basu2006algorithms} for a detailed proof of the structure theorem for subresultants, as well as the explicit formulas for the scaling factors.

\paragraph{The subresultant algorithm.}

The structure theorem for subresultants provides a way to compute the sequence of subresultant polynomials, and hence the GCD and resultant of two polynomials, without performing any division in the fraction field of $R$. Algorithm \ref{alg:subres} gives the pseudo-code of the subresultant algorithm, with Lazard's optimization for computing the scaling factors \cite{ducos2000optimizations}.

\begin{algorithm}
    \caption{\textsf{SubresultantPolynomials}}\label{alg:subres}
    \begin{algorithmic}[1]
        \Require $f,g\in R[x]$, $p=\deg(f)$, $q=\deg(g)$, $p>q$
        \Ensure $(\sResP_i(f,g))_{0\leq i\leq q}$
        \State $\sResP_{p} \gets f$
        \State $s_p \gets 1$, $t_p \gets 1$
        \State $\sResP_{p-1} \gets g$
        \State $t_{p-1} \gets \lc(\sResP_{p-1})$
        \State $\sResP_{q} \gets (-1)^{(p-q)(p-q-1)/2} \lc(\sResP_{p-1})^{p-q-1}\sResP_{p-1}$
        \State $s_q \gets \lc(\sResP_{q})$
        \State $i \gets p+1$, $j \gets p$
        \While{$\sResP_{j-1}\neq 0$}
        \State $k \gets \deg(\sResP_{j-1})$
        \State // Lazard's optimization for computing $s_{k}=\sRes_{k}(f,g)$
        \For{$\delta=1$ to $j-k-1$}
        \State $t_{j-\delta-1} \gets (-1)^{\delta}(t_{j-1}t_{j-\delta})/s_j$
        \EndFor
        \State $s_{k} \gets t_{k}$
        \State // Computation of defective subresultant polynomials
        \State $\sResP_{k} \gets s_k \sResP_{j-1}/t_{j-1}$
        \For{$\ell=k+1$ to $j-1$}
        \State $\sResP_{\ell} \gets 0$
        \EndFor
        \State // Computation of the next non-defective subresultant polynomial
        \State $\sResP_{k-1} \gets -\Rem(s_{k}t_{j-1}\sResP_{i-1},\sResP_{j-1})/(s_j t_{i-1})$
        \State $t_{k-1} \gets \lc(\sResP_{k-1})$
        \State $i \gets j$, $j \gets k$
        \EndWhile
        \For{$\ell=0$ to $j-2$}
        \State $\sResP_{\ell} \gets 0$
        \EndFor
        \\
        \Return $(\sResP_i)_{0\leq i\leq q}$
    \end{algorithmic}
\end{algorithm}

We emphasize that for the subresultant polynomials to be well-defined in $R[x]$, the assumption that $p>q$ is crucial. In the case of $p=q$, the algorithm defines $\sResP_{q}(f,g)=b_q^{-1} g$, which may fail to lie in $R[x]$ if $b_q$ is not a unit. Nevertheless, the polynomials $\sResP_i(f,g)$ for $0\leq i\leq q-1$ remain well-defined in $R[x]$ and can be computed by the algorithm; we do not pursue the case of $p=q$ in this report for simplicity.

Algorithm \ref{alg:subres} returns the sequence of subresultant polynomials $(\sResP_i(f,g))_{0\leq i\leq q}$, which contains all non-defective subresultant polynomials together with some defective ones. If only the non-defective subresultant polynomials are required, the computation of defective subresultant polynomials on Lines 15–19 can be omitted. Similarly, returning only $\sResP_k(f,g)$ computed on Line 16 produces the subresultant polynomial remainder sequence of $f$ and $g$.
